\PassOptionsToPackage{hyphens}{url}
\documentclass[aps,prd,twocolumn,superscriptaddress,nofootinbib]{revtex4-2}

% MiKTeX REVTeX 4.2 needs explicit package loading
\usepackage{amsmath}
\usepackage{amssymb}
\usepackage{amsthm}
\usepackage{booktabs}
\usepackage{tabularx}
\usepackage{xcolor}
\usepackage{graphicx}
\graphicspath{{../figures/}}

\newtheorem{definition}{Definition}
\newtheorem{proposition}{Proposition}
\newtheorem{theorem}{Theorem}
\newtheorem{corollary}{Corollary}
\newtheorem{lemma}{Lemma}
\newtheorem{remark}{Remark}

\usepackage{hyperref}
\hypersetup{
    pdftitle={QG-CRM: Ultraviolet Completion of the Curvature Relaxation Model via Quantum Quadratic Gravity},
    pdfauthor={Lukas Geiger},
    colorlinks=true,
    linkcolor=black,
    urlcolor=blue,
    citecolor=black
}

\begin{document}

% ===================================================================
% TITLE PAGE
% ===================================================================

\title{QG-CRM: Ultraviolet Completion of the Curvature Relaxation Model\\via Quantum Quadratic Gravity}

\author{Lukas Geiger}
\thanks{ORCID: \url{https://orcid.org/0009-0005-7296-1534}}
\affiliation{Independent Researcher, Bernau, Germany}

\date{\today}

\begin{abstract}
Papers~I--IV of the Curvature Relaxation Model (CRM) series demonstrated that late-time cosmic acceleration and galactic dynamics can be explained geometrically via a saturation mechanism $\Omega_\Phi(a) = \Phi_0\tanh(k(a - a_\mathrm{trans}))$, eliminating both dark energy and particulate dark matter. Paper~V elevated this profile to a \textit{universal normal form} via the Saturation Theorem, with the open question: \textit{which ultraviolet completion selects $k$ and $\Phi_0$?} This paper answers that question. We show that the $\gamma R^2$ term in the CRM Lagrangian (Paper~III) admits a natural ultraviolet completion via asymptotically free quantum quadratic gravity (QQG), following the framework of Liu, Quintin \& Afshordi~\cite{LiuQuintinAfshordi2026}. In the UV regime, 1-loop renormalization group running of the $R^2$ coupling $\xi(\mu)$ dynamically generates slow-roll inflation without an inflaton field. We verify that the four axioms of the Saturation Theorem are satisfied by the QQG renormalization group flow, establishing the $\tanh$ profile as the unique infrared attractor. The resulting framework---\textit{QG-CRM}---provides a complete geometric narrative: (i)~asymptotically free $R^2$ gravity drives inflation in the UV; (ii)~saturation dynamics governs the graceful exit and the emergence of general relativity; (iii)~curvature relaxation produces geometric dark matter and late-time acceleration in the IR. We derive the explicit parameter mapping from QQG couplings ($\lambda_0$, $\mathcal{N}$, $\mu_0$) to CRM observables ($k$, $\Phi_0$, $a_\mathrm{trans}$). The tensor-induced inflation scenario of Bertacca et~al.~\cite{Bertacca2025} is classified as a structurally distinct limiting case that falls outside the Saturation Theorem axioms. Observational predictions include $n_s \sim 1 - 4/(3N)$, $r \gtrsim 0.01$, $w(z) < -1$ at late times, and the MOND consistency relation $a_0 = cH_0/(2\pi)$---all testable with Stage~IV CMB experiments, Euclid, and Roman.
\end{abstract}

\keywords{Curvature Relaxation Model, quantum quadratic gravity, ultraviolet completion, inflation, saturation dynamics, dark energy, modified gravity, asymptotic freedom}

\thanks{AI tools (Claude, Anthropic) were used for mathematical formalization and text generation. All physical hypotheses, scientific interpretation, and responsibility for the content lie solely with the author.}

\thanks{Paper~VI in the CRM series. Companion papers: \cite{Geiger2026,Geiger2026b,Geiger2026c,Geiger2026d,Geiger2026e}.}

\maketitle


% ===================================================================
% 1. INTRODUCTION
% ===================================================================
\section{Introduction}
\label{sec:introduction}

The Curvature Relaxation Model (CRM) provides a geometric alternative to the standard $\Lambda$CDM cosmology. Over five companion papers, the framework has been developed from game-theoretic axioms (Paper~I~\cite{Geiger2026}), through a unified treatment of dark energy and MOND-like dark matter (Paper~II~\cite{Geiger2026b}), to microscopic foundations via an $R + \gamma R^2$ Lagrangian with P\"oschl-Teller saturation (Paper~III~\cite{Geiger2026c}), the galactic-cosmological nexus via a vector sector (Paper~IV~\cite{Geiger2026d}), and the Saturation Theorem establishing $\tanh$ as the universal normal form of capacity-limited cooperative relaxation (Paper~V~\cite{Geiger2026e}).

The central phenomenological result---the saturation ODE $dX/da = k(1-X^2)$ and its unique solution $X(a) = \tanh(k(a-a_0))$---has been tested against Pantheon+ supernovae~\cite{Scolnic2022}, Planck CMB observables~\cite{Planck2020}, DESI BAO measurements~\cite{DESI2025}, and SPARC rotation curves (Papers~I--IV), achieving $\Delta\chi^2 = -5.5$ vs.\ $\Lambda$CDM with zero early dark energy and only 6 free parameters.

However, the CRM series has so far remained agnostic about the ultraviolet regime. Paper~V concluded with an explicit open question:
\begin{quote}
``The open question being which UV completion selects $k$ and $\Phi_0$.''
\end{quote}

This paper answers that question. We identify the $\gamma R^2$ sector of the CRM Lagrangian with quantum quadratic gravity (QQG), which is asymptotically free in the ultraviolet and whose 1-loop renormalization group (RG) running dynamically generates slow-roll inflation~\cite{LiuQuintinAfshordi2026}. The Saturation Theorem provides the interface: it constrains the infrared behavior to $\tanh$ regardless of UV details, while the UV completion fixes the microscopic parameters.

The resulting framework, which we call \textit{QG-CRM}, provides a complete geometric narrative from the Big Bang to the present epoch---without an inflaton, without a cosmological constant, and without dark matter particles.

We emphasize that this paper establishes \textit{structural necessity}, not numerical closure. The UV completion determines the existence, functional form, and scaling relations of the CRM parameters; their precise numerical values require nonperturbative UV-IR matching analogous to computing hadron masses from the QCD Lagrangian. What we demonstrate is that QG-CRM is the unique consistent framework connecting asymptotically free quadratic gravity to capacity-limited cosmological relaxation.


% ===================================================================
% 2. THE CRM R² SECTOR
% ===================================================================
\section{The CRM $R^2$ Sector and the Ultraviolet Regime}
\label{sec:r2sector}

\subsection{Recap: The CRM Lagrangian}

Paper~III established the effective CRM Lagrangian~\cite{Geiger2026c}:
\begin{equation}
\mathcal{L}_\mathrm{CRM} = \frac{R}{16\pi G} + \gamma R^2 + \frac{1}{2}(\partial\phi)^2 - V_\mathrm{PT}(\phi),
\label{eq:crm-lagrangian}
\end{equation}
where $\gamma R^2$ generates the geometric ``dark matter'' contribution (the scalaron), and the P\"oschl-Teller potential $V_\mathrm{PT}(\phi)$ provides the saturation dynamics driving late-time acceleration.

In the very early universe, when curvature is Planckian, the Einstein-Hilbert term $R/(16\pi G) \propto M_\mathrm{Pl}^2 R$ becomes subdominant relative to the $R^2$ term. The effective action reduces to
\begin{equation}
S_\mathrm{UV} \approx -\int d^4x\,\sqrt{-g}\,\frac{R^2}{\xi},
\label{eq:pure-r2}
\end{equation}
where $\xi \equiv -1/\gamma$ (following the sign convention of~\cite{LiuQuintinAfshordi2026}; note that $\gamma > 0$ in Paper~III corresponds to $\xi < 0$, i.e., the tachyon-free region that the RG trajectory reaches in the IR). During inflation, $\xi > 0$. This is a scale-invariant theory: any constant-$R$ spacetime---Minkowski, de~Sitter, or anti-de~Sitter---extremizes the action~\cite{Hell2024,Hell2025}. The addition of $R^2$ to the Einstein-Hilbert action is also the starting point of Starobinsky inflation~\cite{Starobinsky1980}, while the full quadratic gravity action including $C^2$ is known to be renormalizable~\cite{Stelle1977}.

\subsection{From scale invariance to inflation}

Scale invariance is broken by quantum effects. The renormalization group running of $\xi(\mu)$ introduces a physical scale into the theory, analogous to the quantum conformal anomaly~\cite{CapperDuff1974,Duff1994}. This is precisely the mechanism that transforms exact de~Sitter into quasi-de~Sitter, generating slow-roll inflation without introducing an inflaton field.

On a homogeneous and isotropic FLRW background, the Weyl tensor $C^2$ vanishes identically. Therefore, the full quadratic gravity action
\begin{equation}
S_\mathrm{QG} = -\int d^4x\,\sqrt{-g}\left(\frac{R^2}{\xi} + \frac{C^2}{2\lambda}\right)
\label{eq:quad-gravity}
\end{equation}
reduces to pure $R^2$ gravity for the cosmological background, and the inflationary dynamics is entirely determined by the running of $\xi$.


% ===================================================================
% 3. UV COMPLETION VIA QQG
% ===================================================================
\section{Ultraviolet Completion via Quantum Quadratic Gravity}
\label{sec:uv-completion}

\subsection{The renormalization group flow}

Following~\cite{Buccio2024,LiuQuintinAfshordi2026}, the 1-loop beta functions of the physical running couplings are
\begin{align}
\beta_\xi &= \frac{d\xi}{d\ln\mu} = -\frac{1}{(4\pi)^2}\,\frac{\xi^2 - 36\lambda\xi - 2520\lambda^2}{36}, \label{eq:beta-xi}\\
\beta_\lambda &= \frac{d\lambda}{d\ln\mu} = -\frac{1}{(4\pi)^2}\,\frac{(1617 + 90\mathcal{N})\lambda - 20\xi}{90}\,\lambda, \label{eq:beta-lambda}
\end{align}
where $\mu$ is the physical RG scale and $\mathcal{N}$ counts the matter field content:
\begin{equation}
\mathcal{N} = \frac{1}{60}N_\mathrm{scalar} + \frac{1}{5}N_\mathrm{vector} + \frac{1}{20}N_\mathrm{fermion}.
\label{eq:matter-content}
\end{equation}

\subsection{Asymptotic freedom and the inflationary trajectory}

The RG flow (\ref{eq:beta-xi})--(\ref{eq:beta-lambda}) admits solutions that are:
\begin{enumerate}
\item \textbf{Asymptotically free in the UV}: both $\xi$ and $\lambda$ vanish at the Gaussian fixed point ($\mu \to \infty$).
\item \textbf{Tachyon-free in the IR}: the low-$\mu$ basin of attraction has $\lambda > 0$ and $\xi < 0$.
\end{enumerate}

Trajectories reaching the tachyon-free region must first pass through a regime where $\xi > 0$: the coupling grows, reaches a maximum $\xi_\mathrm{max}$, then decreases and crosses zero at a scale $\mu_0$. Inflation occurs in the regime of \textit{decreasing} $\xi$, before the tachyon divide is crossed~\cite{LiuQuintinAfshordi2026}.

\subsection{Large-$\mathcal{N}$ solution and the inflation potential}

In the large-$\mathcal{N}$ limit, the solution reads~\cite{LiuQuintinAfshordi2026}:
\begin{equation}
\xi(\mu) \simeq \frac{35\lambda_0^2\ln(\mu/\mu_0)}{8\pi^2[1 + \lambda_\mathrm{tH}\ln(\mu/\mu_0)]},
\label{eq:xi-solution}
\end{equation}
where $\lambda_\mathrm{tH} \equiv \lambda_0\mathcal{N}/(4\pi)^2$ is the 't~Hooft-like coupling. Promoting $\mu = |R|^{1/2}$ and transforming to the Einstein frame yields the scalar potential~\cite{LiuQuintinAfshordi2026}:
\begin{equation}
V(\varphi) \simeq \frac{35\lambda_0^2\mu_0^4}{128\pi^2\lambda_\mathrm{tH}}\left(1 - \frac{\sqrt{6}\,\mu_0}{\lambda_\mathrm{tH}\,\varphi}\right).
\label{eq:einstein-potential}
\end{equation}

This potential has a near-plateau structure, yielding slow-roll inflation with observables~\cite{LiuQuintinAfshordi2026}:
\begin{equation}
n_s \sim 1 - \frac{4}{3N}, \qquad r \sim \frac{8}{3}\left(\frac{2}{\lambda_\mathrm{tH}^2 N^4}\right)^{1/3},
\label{eq:ns-r}
\end{equation}
where $N$ is the number of e-folds between horizon exit and the end of inflation.

\subsection{End of inflation and the onset of kination}

Inflation ends as the potential steepens; there is no minimum where the inflaton could oscillate and decay. Instead, the field rolls into a kinetic-dominated phase (\textit{kination}). As $\lambda$ flows beyond $\lambda_0$, the 't~Hooft coupling $\lambda_\mathrm{tH} = \lambda\mathcal{N}/(4\pi)^2$ surpasses unity, signaling the strong-coupling regime. General relativity must then emerge as the effective field theory at low energies, and the universe enters its radiation era via reheating~\cite{LiuQuintinAfshordi2026}.

This UV$\to$IR transition is precisely where the CRM's infrared dynamics takes over.


% ===================================================================
% 4. THE SATURATION THEOREM AS UV-IR INTERFACE
% ===================================================================
\section{The Saturation Theorem as UV-IR Interface}
\label{sec:interface}

\subsection{Axiom verification for QQG}

Paper~V established four axioms that any quantum gravity theory must satisfy to produce bounded cosmological relaxation~\cite{Geiger2026e}. We now verify that each axiom is satisfied by QQG in the cosmological sector.

\begin{description}
\item[Axiom A (Finite geometric capacity).] The coupling $\xi(\mu)$ reaches a finite maximum $\xi_\mathrm{max}$ along the RG trajectory before decreasing. This maximal curvature response is the geometric capacity bound. \checkmark

\item[Axiom B (Cooperative reinforcement).] The RG running of $\xi$ is self-reinforcing in the inflationary regime: the $\xi^2$ term in $\beta_\xi$ generates cooperative growth from the UV fixed point to $\xi_\mathrm{max}$, followed by cooperative relaxation as $\xi$ decreases. \checkmark

\item[Axiom C (Monotone cosmic control).] The physical RG scale $\mu$ decreases monotonically as the universe expands. During inflation and matter domination, $\mu = |R|^{1/2}$ provides the natural identification. In the radiation era, where the trace $R = 0$ for conformally coupled matter, $\mu$ should be identified with $H$ (the Hubble scale) or an alternative curvature invariant such as the Kretschner scalar $K^{1/4}$; all choices agree during slow-roll inflation and are generically equivalent up to $\mathcal{O}(1)$ factors~\cite{LiuQuintinAfshordi2026}.\footnote{The $\mathcal{O}(1)$ equivalence holds because during quasi-de~Sitter expansion $R \approx 12H^2$, so $|R|^{1/2} \approx 2\sqrt{3}\,H$ and $K^{1/4} \approx 2\sqrt{3}\,H$. In the radiation era, $R = 0$ but $K = 12H^4 > 0$ and $H$ itself remains monotone. The inflationary predictions depend only on the slow-roll regime where all choices coincide; the radiation-era ambiguity affects at most the end-of-inflation matching at the $\mathcal{O}(1)$ level.} The monotonicity of cosmic expansion ensures that $\mu$ decreases in all phases. \checkmark

\item[Axiom D (RG composition).] The 1-loop beta functions (\ref{eq:beta-xi})--(\ref{eq:beta-lambda}) define a well-posed composition law: evolving from $\mu_1 \to \mu_2 \to \mu_3$ is equivalent to evolving directly from $\mu_1 \to \mu_3$. This is precisely the semigroup property of the RG flow, satisfying the Abel equation requirement of the Saturation Theorem. \checkmark
\end{description}

\subsection{Consequence: tanh as the unique IR attractor}

Since all four axioms are satisfied, the Saturation Theorem applies: the unique smooth, scale-coherent response function extending regularly to the saturation boundary is
\begin{equation}
S(g) = S_\mathrm{max}\,\tanh(\alpha\,u(g)),
\end{equation}
up to reparametrization. In cosmological language, this is the CRM saturation profile $\Omega_\Phi(a) = \Phi_0\tanh(k(a - a_\mathrm{trans}))$.

\subsection{Parameter identification}

The UV completion determines the reparametrization $u(g)$ and the CRM constants:
\begin{itemize}
\item The saturation amplitude $\Phi_0$ is fixed by $\xi_\mathrm{max}$ and $\lambda_0$.
\item The relaxation rate $k$ is determined by the slope of $\xi(\mu)$ at the tachyon divide.
\item The transition scale $a_\mathrm{trans}$ corresponds to the scale $\mu_0$ where $\xi = 0$ and GR emerges.
\end{itemize}

The explicit mapping is derived in Appendix~\ref{app:parameter-mapping}.


% ===================================================================
% 5. FROM INFLATION TO LATE-TIME ACCELERATION
% ===================================================================
\section{From Inflation to Late-Time Acceleration: The QG-CRM Narrative}
\label{sec:narrative}

QG-CRM provides a unified geometric narrative spanning the entire cosmic history:

\begin{enumerate}
\item \textbf{Phase~I: Pure $R^2$ inflation (UV, $\mu \gg \mu_0$).} The universe begins in asymptotically free quadratic gravity. The $R^2$ coupling $\xi(\mu)$ runs from zero, reaches $\xi_\mathrm{max}$, and generates quasi-de~Sitter inflation via the logarithmic correction to the R$^2$ action. No inflaton field exists; the scalaron is purely geometric.

\item \textbf{Phase~II: Kination and strong coupling ($\mu \sim \mu_0$).} As $\xi$ decreases and crosses zero, the potential steepens, inflation ends, and the field enters kination. The 't~Hooft coupling approaches unity, signaling the strong-coupling regime. General relativity emerges as the effective field theory.

\item \textbf{Phase~III: Radiation and matter era.} Standard hot Big Bang cosmology. The P\"oschl-Teller scalar $\phi$ is dynamically frozen by the trace coupling $\mathcal{F}(T/\rho)$ (Paper~III), ensuring BBN protection.

\item \textbf{Phase~IV: Geometric dark matter ($z \lesssim 10^3$).} The scalaron from the $\gamma R^2$ term produces scale-dependent gravitational effects that mimic cold dark matter cosmologically (running coupling $\beta_\mathrm{eff}(a)$, Paper~II) and reproduce flat rotation curves via the vector sector (Paper~IV).

\item \textbf{Phase~V: Curvature saturation ($z \lesssim 1$).} The P\"oschl-Teller potential activates, the curvature return potential saturates ($\Omega_\Phi \to \Phi_0$), and the universe enters accelerated expansion. The saturation is guaranteed to follow $\tanh$ by the Saturation Theorem.
\end{enumerate}

This five-phase narrative is not a collection of separate models stitched together. All phases arise from a \textit{single} geometric degree of freedom---the $R^2$ scalaron---whose dynamics is governed by renormalization group flow in the UV and saturation dynamics in the IR, with the Saturation Theorem providing the unique bridge.

\subsection{Scale-dependent scalaron mass: one field, many roles}
\label{sec:scalaron-mass}

A potential objection is that the scalaron cannot simultaneously drive inflation ($m_s \sim 10^{13}$\,GeV in Starobinsky inflation) and produce geometric dark matter effects ($m_s \sim H_0 \sim 10^{-33}$\,eV at cosmological scales). The resolution lies in the \textit{running} of $\gamma(\mu)$: since $\gamma = -1/\xi$ and $\xi(\mu)$ evolves along the RG trajectory, the scalaron mass
\begin{equation}
m_s^2(\mu) = \frac{1}{6\gamma(\mu)} = -\frac{\xi(\mu)}{6}
\label{eq:running-mass}
\end{equation}
is not a fixed parameter but a scale-dependent effective quantity. In the UV ($\mu \gg \mu_0$), $\xi$ is large and positive, yielding a heavy scalaron that drives inflationary dynamics. As $\xi$ decreases through the RG flow, the effective mass drops. After the strong-coupling transition ($\mu \sim \mu_0$), the scalaron enters the tachyon-free regime ($\xi < 0$, $\gamma > 0$), where its mass is controlled by the IR value of $\gamma$---potentially as light as $m_s \sim H_0$.

The ``identity'' of the scalaron across all five phases is therefore not a claim that a single mass parameter governs all epochs, but that a single \textit{geometric degree of freedom} with a \textit{scale-dependent effective description} provides the underlying dynamics. This is directly analogous to the QCD gluon: the same gauge field is perturbative at high energies and confining at low energies, with a scale-dependent effective coupling $\alpha_s(\mu)$ that changes by orders of magnitude.

\subsection{The $R^2 \to$ nonlinear $f(R)$ transition and solar system screening}
\label{sec:screening}

A crucial constraint arises from the companion paper on dark energy~\cite{GeigerDE2026}: the \textit{Convexity-Screening Incompatibility Theorem} establishes that a strictly convex free-energy functional---as naturally produced by linear $R + \gamma R^2$ gravity---admits a unique global minimum with no secondary minima. Consequently, linear $R^2$ gravity \textit{cannot} produce chameleon screening: the post-Newtonian parameter evaluates to $\gamma_\mathrm{PPN} = 1/2$, excluded by Cassini tracking at $>10^4\sigma$.

This has a direct implication for QG-CRM: the UV action (\ref{eq:quad-gravity}) is renormalizably $R^2$, but the IR effective theory \textit{must} develop nonlinear $f(R)$ structure to achieve solar system compatibility. Concretely, the Hu-Sawicki form~\cite{HuSawicki2007}
\begin{equation}
f(R) = R - \mu^2\frac{c_1(R/\mu^2)^n}{c_2(R/\mu^2)^n + 1}, \qquad n \geq 1,
\label{eq:hu-sawicki}
\end{equation}
provides density-dependent screening: the scalaron mass $m_s^2 \propto \rho/M_\mathrm{Pl}^2$ becomes large in high-density environments (solar system), suppressing the Yukawa fifth force, while remaining light ($m_s \sim H_0$) at cosmological densities where dark energy effects are active.

In QG-CRM, this transition has a natural interpretation as an \textit{irrelevant-operator cascade}. In the UV, the theory is perturbatively $R^2$ (renormalizable, with only marginal and relevant operators). As the RG flow approaches strong coupling, nonperturbative effects generate an infinite tower of higher-curvature operators:
\begin{equation}
\Gamma_\mathrm{IR} = \Gamma_{R^2} + \sum_{n \geq 1} \frac{c_n\,R^{n+1}}{\Lambda^{2n}},
\end{equation}
where $\Lambda$ is the confinement scale. The Hu-Sawicki form~(\ref{eq:hu-sawicki}) is then not a fundamental action but the leading \textit{phenomenological resummation} of this operator tower that captures the essential physics: density-dependent scalaron mass and chameleon screening. This is directly analogous to chiral perturbation theory, which parametrizes the low-energy effects of QCD confinement without deriving them from the QCD Lagrangian.

We emphasize that Hu-Sawicki should be understood as an IR effective parametrization of nonperturbative completion, not as a derived action. The specific values of $n$, $c_1$, $c_2$ are constrained by solar system tests and cosmological data, not by the UV theory. What the UV theory guarantees is that \textit{some} nonlinear $f(R)$ structure must emerge at the strong-coupling transition---the Convexity-Screening Theorem makes this a structural necessity, not an optional addition.

\subsection{Connection to the dark energy paper: thermodynamic consistency}
\label{sec:de-connection}

The companion FST dark energy paper~\cite{GeigerDE2026} derives a thermodynamic bound on the effective cosmological constant:
\begin{equation}
\Lambda_\mathrm{eff} \sim \frac{H_0^2}{\ln(M_\mathrm{Pl}/H_0)},
\label{eq:lambda-bound}
\end{equation}
yielding $\Lambda_\mathrm{eff} \approx 3.8 \times 10^{-53}\,\mathrm{m}^{-2}$, within an order of magnitude of the Planck~2018 value ($1.1 \times 10^{-52}\,\mathrm{m}^{-2}$). Since the same Lagrangian ($R + \gamma R^2$ + P\"oschl-Teller) underlies both works, the thermodynamic bound provides an order-of-magnitude \textit{lower bound} on the vacuum energy density. Note that the CRM saturation amplitude $\Phi_0 \equiv \Omega_\Lambda \approx 0.685$ (the fractional dark energy density parameter) is a different quantity from $\Lambda_\mathrm{eff}/(3H_0^2)$. The thermodynamic functional selects the vacuum energy \textit{scale}; the saturation amplitude $\Phi_0$ is determined by the full cosmic expansion history. The ratio
\begin{equation}
\frac{\Lambda_\mathrm{eff}^\mathrm{thermo}}{3H_0^2} \sim \frac{1}{3\ln(M_\mathrm{Pl}/H_0)} \approx 0.0024
\end{equation}
indicates that the raw thermodynamic bound underestimates $\Omega_\Lambda$ by a factor $\sim 290 \approx 2\ln(M_\mathrm{Pl}/H_0)$, identifying the RG matching factor needed to close the gap (see Section~\ref{sec:two-loop}).

Furthermore, the DESI-compatible best-fit parameters from~\cite{GeigerDE2026} ($\kappa = 2.63$, $a_\mathrm{trans} = 0.326$) provide concrete target values for the QQG$\to$CRM parameter mapping (Appendix~\ref{app:parameter-mapping}).

A particularly intriguing prospect is that the QQG renormalization group flow may close the principal open problem identified in~\cite{GeigerDE2026}: the missing factor $\sim 2\ln(M_\mathrm{Pl}/H_0) \approx 290$. The thermodynamic functional provides one logarithm from the entropy/state-density logic; a second logarithm arises naturally from the RG integration of the vacuum energy running:
\begin{equation}
\Lambda_\mathrm{eff} \sim \int_{H_0}^{M_\mathrm{Pl}} \frac{d\ln\mu}{\mu}\,\beta_\Lambda(\mu) \sim \frac{H_0^2\,\ln^2(M_\mathrm{Pl}/H_0)}{(4\pi)^2},
\end{equation}
where the double-logarithm structure is a standard feature of integrated RG flows in quantum field theory. This ``double-log'' mechanism---one $\ln$ from thermodynamics, one from UV$\to$IR integration---would yield $\Phi_0 \sim \ln(M_\mathrm{Pl}/H_0)/(3\ln(M_\mathrm{Pl}/H_0)) \times \mathcal{O}(\text{few}) \sim \mathcal{O}(1)$, consistent with $\Omega_\Lambda \approx 0.685$. A rigorous derivation requires explicit evaluation of $\beta_\Lambda$ in QQG, which we defer to future work.


% ===================================================================
% 6. RELATION TO TENSOR-INDUCED INFLATION
% ===================================================================
\section{Relation to Tensor-Induced Inflation}
\label{sec:tensor-inflation}

Bertacca et~al.~\cite{Bertacca2025} proposed a scenario---building on the foundational work of Matarrese, Mollerach \& Bruni~\cite{Matarrese1998}---in which scalar perturbations arise purely as a second-order effect from tensor fluctuations in exact de~Sitter, with no scalar degrees of freedom at linear order. While this is an elegant construction, it is structurally incompatible with saturation cosmology for the following reasons:

\begin{enumerate}
\item \textbf{No linear scalar sector.} The scenario explicitly excludes linear scalar perturbations. The CRM, by contrast, requires the scalaron as a dynamical degree of freedom for both the inflationary background and the late-time saturation mechanism.

\item \textbf{No saturation dynamics.} Without a scalar order parameter, there is no composable relaxation variable satisfying the Abel equation. The Saturation Theorem is therefore not applicable in this setting, and the $\tanh$ normal form cannot be derived from the tensor-only framework.

\item \textbf{No global cosmological program.} The tensor-induced scenario addresses inflation and perturbation generation, but does not extend to dark matter, MOND, or late-time acceleration.
\end{enumerate}

However, the tensor-to-scalar sourcing mechanism of~\cite{Bertacca2025} may provide a \textit{subleading correction} to the scalaron-dominated power spectrum in QG-CRM. The quantitative assessment of this correction---whether the second-order tensor channel modifies $\mathcal{P}_\zeta(k)$ at the percent level---remains an open question for future work.

\begin{remark}
The scenario of~\cite{Bertacca2025} predicts a very different tensor-to-scalar ratio ($r \sim 0.0006$, near GUT scale $H_\mathrm{inf} \sim 5 \times 10^{-4}\,m_\mathrm{pl}$) from the QQG prediction ($r \gtrsim 0.01$). Stage~IV CMB experiments will decisively distinguish these scenarios.
\end{remark}


% ===================================================================
% 7. OBSERVATIONAL PREDICTIONS
% ===================================================================
\section{Observational Predictions of QG-CRM}
\label{sec:predictions}

QG-CRM makes concrete, falsifiable predictions spanning the entire cosmic history:

\subsection{Inflationary epoch (UV)}

We have computed the inflationary observables numerically across the viable parameter space. The results are summarized in Table~\ref{tab:observables}.

\begin{table}[h]
\caption{QG-CRM inflationary predictions for representative parameter values. The observed scalar amplitude $A_s = e^{3.06} \times 10^{-10}$ is used to fix $\lambda_0$ for each $(\lambda_\mathrm{tH}, N)$ pair.}
\begin{tabular}{cccccc}
\toprule
$\lambda_\mathrm{tH}$ & $N$ & $n_s$ & $r$ & $\lambda_0$ & $\mathcal{N}$ \\
\midrule
0.3 & 55 & 0.9758 & 0.036 & $1.1 \times 10^{-4}$ & $4.3 \times 10^5$ \\
0.5 & 55 & 0.9758 & 0.025 & $1.2 \times 10^{-4}$ & $6.6 \times 10^5$ \\
1.0 & 55 & 0.9758 & 0.016 & $1.4 \times 10^{-4}$ & $1.2 \times 10^6$ \\
\midrule
\multicolumn{2}{c}{Starobinsky ($N=55$)} & 0.9636 & 0.004 & --- & --- \\
\multicolumn{2}{c}{ACT+DESI (1$\sigma$)} & $0.974 \pm 0.003$ & --- & --- & --- \\
\multicolumn{2}{c}{SPT-3G+DESI (1$\sigma$)} & $0.978 \pm 0.004$ & --- & --- & --- \\
\multicolumn{2}{c}{BICEP/Keck (2$\sigma$)} & --- & $<0.036$ & --- & --- \\
\bottomrule
\end{tabular}
\label{tab:observables}
\end{table}

\begin{itemize}
\item \textbf{Scalar spectral index:} $n_s = 1 - 4/(3N) = 0.9758$ for $N = 55$, lying within $1\sigma$ of both ACT+DESI~\cite{ACT2025} and SPT-3G+DESI~\cite{SPT3G2026}. This is significantly larger than Starobinsky's $n_s = 1 - 2/N = 0.964$, which is under tension with recent CMB data.

\item \textbf{Tensor-to-scalar ratio:} $r$ ranges from $0.016$ ($\lambda_\mathrm{tH} = 1$, at the strong-coupling boundary) to $0.036$ ($\lambda_\mathrm{tH} = 0.3$), with the minimum prediction $r \gtrsim 0.01$. The upper end is at the current BICEP/Keck $2\sigma$ bound. LiteBIRD ($\sigma_r \sim 0.001$) and CMB-S4 will probe the entire viable range.

\item \textbf{Scalar amplitude:} Fixing $A_s = 2.13 \times 10^{-9}$ constrains $\lambda_0 \sim 10^{-4}$ and requires $\mathcal{N} \sim 10^5$--$10^6$ matter fields in the loops. This large $\mathcal{N}$ is compatible with extensions of the standard model and with holographic cosmology~\cite{Afshordi2017}.

\item \textbf{Tensor-sourced correction.} The second-order tensor$\to$scalar contribution of~\cite{Bertacca2025} is suppressed by $(H_\mathrm{inf}/m_\mathrm{pl})^2 \sim 2.5 \times 10^{-7}$ relative to the scalaron channel, confirming its subleading nature.
\end{itemize}

\subsection{Late-time cosmology (IR)}

Inherited from Papers~I--IV:
\begin{itemize}
\item \textbf{Equation of state and phantom stability:} $w_\mathrm{eff}(z) < -1$ with $|\Delta w| \approx 0.4$ (phantom-like). Crucially, this phantom phase is \textit{transient}, not asymptotic. As $a \to \infty$: $\Omega_m a^{-3} \to 0$, $\Omega_\Phi \to \Phi_0$, hence $H^2 \to H_0^2\Phi_0 = \mathrm{const}$ and $\dot{H} \to 0$. The universe approaches a stable de~Sitter end state with $w_\mathrm{eff} \to -1$ from below. No Big Rip, no future singularity---the saturation bound $\Omega_\Phi \leq \Phi_0$ is the structural guarantee.
\item \textbf{Later acceleration onset:} $z_\mathrm{acc} = 0.52$ vs.\ $0.84$ ($\Lambda$CDM). The extended matter-dominated phase explains JWST ``Universe Breakers'' at $z > 7$~\cite{Labbe2023}.
\item \textbf{Modified gravity signatures:} $\mu(k,a) \neq 1$, gravitational slip $\eta \neq 1$, lensing $\Sigma = 1$---distinguishing QG-CRM from $\Lambda$CDM. Testable with Euclid and Roman.
\item \textbf{Gravitational wave speed:} $c_T = c$ (exact), consistent with GW170817~\cite{GW170817}.
\item \textbf{MOND consistency relation:} $a_0 = cH_0/(2\pi)$~\cite{Milgrom1983}, linking local galactic dynamics to the cosmological expansion rate.
\end{itemize}


% ===================================================================
% 8. DISCUSSION
% ===================================================================
\section{Discussion}
\label{sec:discussion}

\subsection{The ontological picture}

QG-CRM reduces the material content of cosmology to its minimum:
\begin{center}
\textit{Spacetime geometry (in five phases) + baryonic matter.}
\end{center}
No cosmological constant, no inflaton field, no dark matter particles, no dark energy fluid. The entire cosmic history---from the Big Bang singularity to the present accelerated expansion---is governed by a single geometric degree of freedom (the $R^2$ scalaron) whose behavior transitions from UV renormalization group dynamics to IR saturation dynamics.

\subsection{Implications for Paper~I: The Nash equilibrium in the ultraviolet}
\label{sec:paper1-implications}

The UV completion has direct consequences for the game-theoretic framework of Paper~I~\cite{Geiger2026}, which modeled cosmic evolution as a Nash equilibrium between a metastable quantum vacuum (``null space'') and a spacetime bubble. We can now assign physical identities to both players.

\paragraph{Player identification.}
The null space of Paper~I corresponds to the pure $R^2$ regime of QQG: a scale-invariant theory without classical spacetime, where the universe exists as a quantum state near the Gaussian UV fixed point. The spacetime bubble corresponds to the emergent GR regime after the strong-coupling transition at $\mu \sim \mu_0$. The Nash equilibrium---the endpoint of the game---maps to the $\tanh$ saturation attractor, which the Saturation Theorem guarantees as the unique stable outcome.

\paragraph{The complete five-phase game.}
Paper~I described only the late-time phases (geometric dark matter $\to$ saturation). QG-CRM extends the game to all five phases:
\begin{enumerate}
\item \textit{Inflation}: the null-space player dominates (pure $R^2$, no GR).
\item \textit{Kination/reheating}: the ``move'' passes from player~1 to player~2 (strong-coupling transition, GR emerges).
\item \textit{Radiation/matter}: the spacetime-bubble player dominates (standard cosmology).
\item \textit{Geometric dark matter}: cooperative phase (both players contribute).
\item \textit{Saturation}: Nash equilibrium is reached (de~Sitter end state).
\end{enumerate}

\paragraph{Parameter derivation: what the UV fixes and what it does not.}
In Paper~I, the saturation amplitude $\Phi_0$ and relaxation rate $k$ were fitted to Pantheon+ data. The UV completion clarifies the \textit{logical status} of these parameters: QQG determines (i)~the \textit{existence} of a finite saturation value (via $\xi_\mathrm{max}$), (ii)~the \textit{functional form} of relaxation ($\tanh$, guaranteed by the Saturation Theorem), and (iii)~the \textit{initial conditions} for the saturation ODE at the reheating surface. However, the \textit{numerical values} of $\Phi_0$ and $k$ are IR quantities, determined jointly by the UV initial conditions and the cosmological boundary conditions ($H_0$, $\Omega_m$). The strong-coupling transition is the correct UV-IR matching surface, but it transfers the dynamics to the IR regime rather than fixing $\Phi_0$ directly. A fully quantitative derivation would require tracking the RG flow through the nonperturbative strong-coupling window---an open problem analogous to computing hadron masses from the QCD Lagrangian.

\paragraph{Potential identification: projection, not identity.}
The game-theoretic potential of Paper~I and the QQG Einstein-frame potential~(\ref{eq:einstein-potential}) are not the same object, but represent \textit{different levels of description of the same system}. The QQG potential is a microscopic RG-improved potential, valid in the UV/inflationary regime. The Paper~I potential is an effective IR free-energy potential, defined through relaxation dynamics and thermodynamic selection. The connection between them is mediated by the Saturation Theorem: the QQG potential provides the UV initial conditions, and the $\tanh$ normal form projects these onto the unique IR attractor. In Lagrangian language, the Paper~I potential corresponds to the effective potential of the scalaron in the Einstein frame, evaluated at cosmological densities and integrated over the full RG flow including the strong-coupling transition.

\paragraph{Thermodynamic strengthening.}
Paper~I established a thermodynamic equivalence connecting the game-theoretic axioms to the Jacobson tradition~\cite{Jacobson1995}. The UV completion strengthens this connection through three independent results: (i)~the quantum-thermodynamic de~Sitter decay mechanism~\cite{Alicki2023,Alicki2025}, (ii)~Mottola's thermodynamic instability of de~Sitter space~\cite{Mottola1985,Mottola1986}, and (iii)~the quantum break-time of de~Sitter~\cite{DvaliGomezZell2017}. All three provide concrete physical mechanisms for the ``game dynamics'' that Paper~I described axiomatically.

\subsection{Comparison with alternative frameworks}

\begin{table}[h]
\caption{Comparison of cosmological frameworks.}
\begin{tabular}{lccc}
\toprule
 & $\Lambda$CDM & Starobinsky & QG-CRM \\
\midrule
Inflaton field & Yes & Yes (scalaron) & No \\
Cosmological constant & Yes & No & No \\
Dark matter particles & Yes & Yes & No \\
UV complete & No & No & Yes$^\dagger$ \\
$n_s$ prediction & --- & $1-2/N$ & $1-4/(3N)$ \\
$r$ prediction & --- & $12/N^2$ & $\gtrsim 0.01$ \\
MOND emergence & No & No & Yes \\
\bottomrule
\end{tabular}
\label{tab:comparison}
\end{table}
\noindent$^\dagger$\,Conditional on ghost confinement in the strong-coupling regime (see Sec.~\ref{sec:ghost}).

\subsection{Robustness of the large-$\mathcal{N}$ requirement}
\label{sec:large-n}

The viable parameter space requires $\mathcal{N} \sim 10^5$--$10^6$ matter fields contributing to the QQG beta functions. We identify three classes of beyond-standard-model (BSM) scenarios that naturally provide this:

\begin{enumerate}
\item \textbf{Large gauge groups.} An $\mathrm{SU}(N_c)$ theory with $N_c \sim 300$--$1000$ and a small number of fundamental fermions yields $\mathcal{N} \sim N_c^2/5 \sim 10^4$--$2 \times 10^5$ from the adjoint gauge bosons alone. This is the scenario most closely related to holographic cosmology~\cite{Afshordi2017}, where $N_c^2$ plays the role of the central charge.

\item \textbf{Large hidden sectors.} String landscape constructions with $\sim 10^5$ SM-like gauge sectors yield $\mathcal{N} \sim 5 \times 10^5$ while preserving individual sector simplicity.

\item \textbf{Scalar moduli.} Type~IIB flux compactifications with $\sim 6 \times 10^6$ scalar moduli give $\mathcal{N} \sim 10^5$ from the scalar contribution alone ($\mathcal{N}_\mathrm{scalar} = N_\mathrm{scalar}/60$).
\end{enumerate}

We stress the distinction between \textit{existential} and \textit{realistic} large $\mathcal{N}$: the QG-CRM framework \textit{requires} $\mathcal{N} \sim 10^5$--$10^6$ for perturbative control, but it does not specify which fields provide this. The scenarios above demonstrate that such field content is \textit{available} in well-motivated BSM settings, not that it is uniquely determined.

Crucially, these fields do \textit{not} contribute to the effective number of relativistic species $N_\mathrm{eff}$ at BBN or CMB decoupling. This is a structural feature, not a fine-tuning: in QG-CRM, the large-$\mathcal{N}$ fields are active \textit{exclusively} above the strong-coupling scale ($\mu > \mu_0$). At the confinement transition, they are integrated out together with the QQG degrees of freedom---precisely as heavy quarks are integrated out above their mass threshold in QCD. Below the confinement scale, only standard model fields ($\mathcal{N}_\mathrm{SM} \approx 4.7$) propagate as asymptotic states, yielding $\Delta N_\mathrm{eff} = 0$ from the hidden sector. Entropy production at reheating is $S_f/S_i \sim 1$--$3$, consistent with standard baryogenesis.

\subsection{Ghost confinement in the Weyl sector}
\label{sec:ghost}

The $C^2$ term in quadratic gravity introduces a massive spin-2 field with negative kinetic energy---a ``ghost'' in the perturbative spectrum. However, this does not imply the existence of propagating ghost \textit{asymptotic states}. In QG-CRM, the ghost is a \textit{confined resonance}, not an on-shell particle: as $\lambda$ grows toward the strong-coupling regime ($\lambda_\mathrm{tH} \gtrsim 1$), the ghost becomes confined, in direct analogy to how gluons in QCD appear in the perturbative spectrum but never as free asymptotic states~\cite{HoldomRen2016}. Three aspects support this:

\begin{enumerate}
\item \textbf{Asymptotic freedom.} Both QQG and QCD are asymptotically free, with perturbative UV behavior and nonperturbative IR dynamics. In QCD, confinement prevents free quarks from appearing as asymptotic states; in QQG, confinement is expected to prevent the ghost from appearing as a physical excitation.

\item \textbf{Dynamical mass generation.} In the strong-coupling regime, the ghost may acquire a dynamical mass $m_\mathrm{ghost} \sim \mu_0$, pushing its effects above the energy scale of cosmological observations. This is analogous to constituent quark masses in QCD.

\item \textbf{Observable signatures.} If ghost confinement is incomplete, the most likely signature would be a modification of the tensor power spectrum at scales $k \sim m_\mathrm{ghost}$, producing a feature in the $B$-mode polarization at high $\ell$. The absence of such features in current CMB data provides an indirect consistency check, though dedicated searches with Stage~IV experiments would be valuable.
\end{enumerate}

We acknowledge that a rigorous proof of ghost confinement in QQG remains outstanding. Lattice formulations of quadratic gravity~\cite{HoldomRen2016} offer the most promising path toward a nonperturbative verification.

\subsection{Nonperturbative scenarios for $\Phi_0$ and $k$}
\label{sec:nonperturbative}

While the precise numerical values of the CRM parameters $\Phi_0$ and $k$ require nonperturbative UV-IR matching, we can sketch two plausible scenarios:

\paragraph{Scenario A: Dimensional transmutation.}
By analogy with $\Lambda_\mathrm{QCD}$, the strong-coupling transition in QQG generates a dynamical scale $\mu_\mathrm{conf} \sim \mu_0\,e^{-1/\lambda_\mathrm{tH}}$. The ratio $\mu_\mathrm{conf}/M_\mathrm{Pl}$ then sets the vacuum energy scale, and the CRM parameters inherit their values from this dynamically generated hierarchy. In this scenario, $\Phi_0 \propto (\mu_\mathrm{conf}/M_\mathrm{Pl})^4$ and $k$ is set by the logarithmic derivative of $\xi$ at the confinement scale.

\paragraph{Scenario B: Thermodynamic selection.}
The FST dark energy paper~\cite{GeigerDE2026} derives $\Lambda_\mathrm{eff} \sim H_0^2/\ln(M_\mathrm{Pl}/H_0)$ from a thermodynamic variational principle. If the QQG RG flow provides the microscopic underpinning for this variational principle---with the free-energy functional $\Phi(\rho_\mathrm{vac})$ emerging from the path integral of QQG in the strong-coupling regime---then $\Phi_0$ is selected by thermodynamic equilibrium rather than by UV dynamics alone.

Both scenarios predict $\Phi_0 \sim \mathcal{O}(1)$ (in units of the critical density), consistent with the observed $\Omega_\Lambda \approx 0.685$, but a derivation from first principles remains a major open challenge.

\subsection{Perturbative stability: two-loop estimate}
\label{sec:two-loop}

We estimate the impact of 2-loop corrections on the inflationary predictions using the parametric scaling $\delta(\mathrm{obs})/\mathrm{obs} \sim \lambda_\mathrm{tH}/(4\pi)^2$. For $\lambda_\mathrm{tH} = 1$ (the strong-coupling boundary), this gives $\delta \sim 0.6\%$. Concretely:

\begin{itemize}
\item The spectral index shift is $|\delta n_s| < 1.5 \times 10^{-4}$, which is 20$\times$ smaller than the current $1\sigma$ observational error ($\sigma_{n_s} \approx 0.003$). The $n_s$ prediction is \textit{robust}.
\item The tensor-to-scalar ratio shift is $|\delta r/r| < 0.6\%$, corresponding to $|\delta r| < 10^{-4}$. This is below LiteBIRD's forecast sensitivity ($\sigma_r \sim 10^{-3}$), though marginally detectable at the $\sim 10\%$-of-$\sigma$ level.
\item Perturbative control breaks down only at $\lambda_\mathrm{tH} \gtrsim 2$, well outside the viable parameter space ($0.1 \lesssim \lambda_\mathrm{tH} \lesssim 1$).
\end{itemize}

We conclude that the 1-loop predictions of Table~\ref{tab:observables} are stable against higher-order corrections throughout the viable parameter space.

\subsection{Remaining open questions}

\begin{enumerate}
\item \textbf{R$^2$ $\to$ nonlinear $f(R)$ transition} [\textit{requires nonperturbative methods}]. The Convexity-Screening Theorem mandates nonlinear $f(R)$ in the IR, but the confinement mechanism that generates Hu-Sawicki structure from the UV $R^2$ action is inherently nonperturbative. Lattice methods for quadratic gravity may provide a path forward.

\item \textbf{RG matching for $\Lambda_\mathrm{eff}$} [\textit{plausible, open}]. The missing factor $\sim 290 \approx 2\ln(M_\mathrm{Pl}/H_0)$ in the thermodynamic bound is suggestive of a 2-loop $\ln^2$ contribution or large-$\mathcal{N}$ enhancement, but the calculation remains open.

\item \textbf{Perturbations from the Weyl sector.} Running of the $C^2$ coupling may modify the tensor power spectrum and hence $r$; a quantitative assessment is needed.
\end{enumerate}

These are genuine \textit{research questions}, not structural deficiencies. The QG-CRM framework is internally consistent at the level presented; the open questions concern quantitative refinements within a fixed architecture.


% ===================================================================
% 9. CONCLUSION
% ===================================================================
\section{Conclusion}
\label{sec:conclusion}

We have presented QG-CRM, the ultraviolet-completed Curvature Relaxation Model. By identifying the $R^2$ sector of the CRM Lagrangian with asymptotically free quantum quadratic gravity, and by verifying that the Saturation Theorem provides the unique UV-IR interface, we obtain a complete geometric cosmology spanning from the Big Bang to the present epoch.

The key results are:

\begin{enumerate}
\item The UV completion is provided by quantum quadratic gravity: renormalization group running of the $R^2$ coupling dynamically generates inflation without an inflaton field.

\item The Saturation Theorem (Paper~V) applies to the QQG RG flow, guaranteeing $\tanh$ as the unique IR normal form and establishing the CRM saturation profile as structurally necessary.

\item The UV completion answers the open question of Paper~V by identifying QQG as the theory that selects the CRM parameters: it determines the existence and functional form of saturation, fixes the inflationary observables, and provides scaling relations for the IR parameters $\Phi_0$ and $k$---though their precise numerical values require nonperturbative UV-IR matching.

\item QG-CRM makes falsifiable predictions---$n_s \sim 1 - 4/(3N)$, $r \gtrsim 0.01$, phantom-like $w(z)$, and $a_0 = cH_0/(2\pi)$---that are testable with Stage~IV CMB experiments, gravitational wave detectors, and galaxy surveys.
\end{enumerate}

QG-CRM is, to our knowledge, the first framework that unifies ultraviolet-complete inflation, geometric dark matter, MOND dynamics, and late-time cosmic acceleration within a single action principle---using only spacetime curvature and baryonic matter.


% ===================================================================
% ACKNOWLEDGMENTS
% ===================================================================
\begin{acknowledgments}
AI tools (Claude, Anthropic; Copilot, Microsoft) were used for mathematical formalization, literature analysis, and text generation. All physical hypotheses, scientific interpretation, and responsibility for the content lie solely with the author.
\end{acknowledgments}


% ===================================================================
% APPENDICES
% ===================================================================
\appendix

\section{Explicit Parameter Mapping: QQG $\to$ CRM}
\label{app:parameter-mapping}

The parameter mapping must be understood at two levels: what the UV completion \textit{determines structurally} and what it \textit{leaves to IR boundary conditions}.

\paragraph{What QQG determines.}
\begin{enumerate}
\item \textbf{Existence of saturation.} The RG trajectory has a finite $\xi_\mathrm{max}$, guaranteeing that the relaxation variable is bounded. This is Axiom~A of the Saturation Theorem.
\item \textbf{Functional form.} The Saturation Theorem forces $\tanh$ as the unique normal form. No alternative profile is permitted.
\item \textbf{Initial conditions.} At the reheating surface ($\mu \sim \mu_0$, where strong coupling sets in and GR emerges), the QQG dynamics sets the initial value and velocity of the relaxation variable.
\item \textbf{Inflationary observables.} $n_s$, $r$, and $A_s$ are fully determined by $\lambda_\mathrm{tH}$ and $N$ (Table~\ref{tab:observables}).
\end{enumerate}

\paragraph{What requires IR boundary conditions.}
The saturation amplitude $\Phi_0$ and relaxation rate $k$ are IR quantities. Their numerical values are determined jointly by:
\begin{itemize}
\item the UV initial conditions at the reheating surface,
\item the cosmological expansion history through radiation and matter domination,
\item the present-day boundary conditions ($H_0$, $\Omega_m$).
\end{itemize}

\paragraph{UV-derived scaling relations.}
While $\Phi_0$ and $k$ cannot be computed from UV parameters alone, the UV completion provides \textit{scaling relations}:

\begin{itemize}
\item The inflation energy scale sets an upper bound:
\begin{equation}
\frac{V_\mathrm{inf}}{M_\mathrm{Pl}^4} = \frac{35\lambda_0^2\mu_0^4}{128\pi^2\lambda_\mathrm{tH}\,M_\mathrm{Pl}^4} \sim 10^{-9},
\end{equation}
confirming that the inflationary vacuum energy is far above the present dark energy scale.

\item The RG slope at the tachyon divide sets the \textit{rate} at which the scalaron coupling evolves:
\begin{equation}
\left.\frac{d\xi}{d\ln\mu}\right|_{\mu=\mu_0} = \frac{2520\lambda_0^2}{36(4\pi)^2} \sim 4.4 \times 10^{-5},
\end{equation}
providing a characteristic timescale for the UV$\to$IR transition.

\item \textbf{The relaxation rate $k$ as a critical exponent.} At the tachyon divide ($\xi = 0$), the linearized RG flow of $\xi$ has the Jacobian
\begin{equation}
\left.\frac{\partial\beta_\xi}{\partial\xi}\right|_{\xi=0} = \frac{\lambda_0}{(4\pi)^2}.
\end{equation}
This quantity plays the role of a \textit{critical exponent}: it characterizes the rate at which the system departs from the critical surface $\xi = 0$, independently of the nonperturbative dynamics beyond this surface. The CRM relaxation rate $k$ inherits this universal scaling:
\begin{equation}
k \propto \frac{\lambda_0}{(4\pi)^2} \sim 7 \times 10^{-7},
\label{eq:k-critical}
\end{equation}
where the numerical value uses $\lambda_0 \sim 10^{-4}$ from the $A_s$ constraint. This is a structural result: $k$ is determined by the \textit{local} properties of the RG flow at the critical surface, not by the global nonperturbative dynamics.

\textbf{Important clarification:} The RG exponent $k_\mathrm{RG} \sim 7 \times 10^{-7}$ and the observed CRM relaxation rate $\kappa = 2.63$ are \textit{not the same quantity}. They differ by the Jacobian of the coordinate transformation between RG time ($\ln\mu$) and cosmic time ($\ln a$):
\begin{equation}
\kappa = k_\mathrm{RG} \cdot \left|\frac{d\ln\mu}{d\ln a}\right| \cdot \mathcal{J}_\mathrm{repar},
\label{eq:k-jacobian}
\end{equation}
where $\mathcal{J}_\mathrm{repar}$ is the reparametrization Jacobian from the Saturation Theorem (which guarantees $\tanh$ ``up to reparametrization''). Since $\mu \sim |R|^{1/2}$ and $R \propto H^2 \propto a^{-3}$ during matter domination, $|d\ln\mu/d\ln a| = 3/2$ per e-fold of expansion. Over the full cosmic evolution from reheating ($a \sim 10^{-28}$) to present ($a = 1$), the integrated Jacobian spans $\sim 64$ e-folds of expansion, each contributing $\mathcal{O}(1)$ to the mapping. The remaining factor $\sim 10^{4\text{--}5}$ arises from the reparametrization between the RG control variable and the CRM relaxation variable---a nontrivial but structurally expected consequence of projecting a UV flow onto an IR normal form.

\item The reheating temperature from out-of-equilibrium dS decay~\cite{LiuQuintinAfshordi2026}:
\begin{equation}
T_\mathrm{RH} = 10^{-10}\left(\frac{E_\mathrm{inf}}{\mathrm{GeV}}\right)^{3/2}\,\mathrm{GeV} \sim 10^{11}\text{--}10^{12}\,\mathrm{GeV},
\end{equation}
consistent with baryogenesis.
\end{itemize}

\paragraph{The nonperturbative gap.}
A fully quantitative derivation of $\Phi_0$ and $k$ from UV parameters would require solving the RG flow through the strong-coupling regime at $\lambda_\mathrm{tH} \sim 1$. This is analogous to computing hadron masses from the QCD Lagrangian: the theoretical framework is established, but the nonperturbative computation remains an open challenge. Lattice methods for quadratic gravity~\cite{HoldomRen2016} may ultimately provide the necessary bridge.

\section{Scalar Perturbation Power Spectrum}
\label{app:power-spectrum}

The primordial scalar power spectrum in QG-CRM has two contributions:

\begin{enumerate}
\item \textbf{Primary (scalaron):} The standard $f(R)$-inflation result from the Einstein-frame potential (\ref{eq:einstein-potential}), yielding~\cite{LiuQuintinAfshordi2026}
\begin{equation}
\Delta_\phi^2(k) \propto \frac{\lambda_0^2 N^{4/3}}{(2\lambda_\mathrm{tH})^{1/3}}.
\end{equation}

\item \textbf{Secondary (tensor-sourced):} The second-order contribution from tensor perturbations via the kernel computed in~\cite{Bertacca2025}. This is suppressed relative to the scalaron channel by a factor of order $(H_\mathrm{inf}/m_\mathrm{pl})^2 \ll 1$ and is therefore subleading.
\end{enumerate}

The detailed computation of the combined spectrum, including interference effects and non-Gaussian corrections, is deferred to future work.


% ===================================================================
% REFERENCES
% ===================================================================
\begin{thebibliography}{99}

\bibitem{Geiger2026}
L.~Geiger,
``Game-Theoretic Cosmology and the Curvature Relaxation Model: A Nash Equilibrium Framework for Cosmic Acceleration'' (Paper~I),
Zenodo (2026), \href{https://doi.org/10.5281/zenodo.18728935}{10.5281/zenodo.18728935} (combined CRM~I--IV record).

\bibitem{Geiger2026b}
L.~Geiger,
``Eliminating the Dark Sector: Unifying the Curvature Relaxation Model with MOND'' (Paper~II),
in Ref.~\cite{Geiger2026}.

\bibitem{Geiger2026c}
L.~Geiger,
``Microscopic Foundations of the Curvature Relaxation Model: From Quantum Geometry to Macroscopic Saturation'' (Paper~III),
in Ref.~\cite{Geiger2026}.

\bibitem{Geiger2026d}
L.~Geiger,
``The Galactic--Cosmological Nexus: Connecting CRM to MOND Dynamics via Curvature Saturation'' (Paper~IV),
in Ref.~\cite{Geiger2026}.

\bibitem{Geiger2026e}
L.~Geiger,
``The Saturation Theorem: Axiomatic Constraints on Quantum Gravity from Cosmological Relaxation,''
Zenodo (2026), \href{https://doi.org/10.5281/zenodo.19036188}{10.5281/zenodo.19036188}.
% CRM Paper V

\bibitem{LiuQuintinAfshordi2026}
R.~Liu, J.~Quintin and N.~Afshordi,
``Ultraviolet Completion of the Big Bang in Quadratic Gravity,''
Phys.\ Rev.\ Lett.\ \textbf{136}, 111501 (2026),
\href{https://doi.org/10.1103/6gtx-j455}{10.1103/6gtx-j455}.

\bibitem{Bertacca2025}
D.~Bertacca, R.~Jimenez, S.~Matarrese and A.~Ricciardone,
``Inflation without an inflaton,''
Phys.\ Rev.\ Research \textbf{7}, L032010 (2025),
\href{https://doi.org/10.1103/vfny-pgc2}{10.1103/vfny-pgc2}.

\bibitem{Buccio2024}
D.~Buccio, J.~F.~Donoghue, G.~Menezes and R.~Percacci,
``Physical running of couplings in quadratic gravity,''
Phys.\ Rev.\ Lett.\ \textbf{133}, 021604 (2024).

\bibitem{Starobinsky1980}
A.~A.~Starobinsky,
``A new type of isotropic cosmological models without singularity,''
Phys.\ Lett.\ B \textbf{91}, 99 (1980).

\bibitem{Hell2024}
A.~Hell, D.~Lust and G.~Zoupanos,
``On the degrees of freedom of $R^2$ gravity in flat spacetime,''
JHEP \textbf{02}, 039 (2024).

\bibitem{Hell2025}
A.~Hell and D.~Lust,
``Conformal and pure scale-invariant gravities in $d$ dimensions,''
JHEP \textbf{09}, 202 (2025).

\bibitem{CapperDuff1974}
D.~M.~Capper and M.~J.~Duff,
``Trace anomalies in dimensional regularization,''
Nuovo Cimento Soc.\ Ital.\ Fis.\ \textbf{23A}, 173 (1974).

\bibitem{Duff1994}
M.~J.~Duff,
``Twenty years of the Weyl anomaly,''
Class.\ Quant.\ Grav.\ \textbf{11}, 1387 (1994).

\bibitem{HoldomRen2016}
B.~Holdom and J.~Ren,
``QCD analogy for quantum gravity,''
Phys.\ Rev.\ D \textbf{93}, 124030 (2016).


\bibitem{GeigerDE2026}
L.~Geiger,
``Dark Energy as Residual Vacuum Free Energy: A Thermodynamic Bound on the Cosmological Constant,''
Zenodo (2026), \href{https://doi.org/10.5281/zenodo.19106160}{10.5281/zenodo.19106160}.

\bibitem{HuSawicki2007}
W.~Hu and I.~Sawicki,
``Models of $f(R)$ cosmic acceleration that evade solar-system tests,''
Phys.\ Rev.\ D \textbf{76}, 064004 (2007).

\bibitem{Stelle1977}
K.~S.~Stelle,
``Renormalization of higher-derivative quantum gravity,''
Phys.\ Rev.\ D \textbf{16}, 953 (1977).

\bibitem{DvaliGomezZell2017}
G.~Dvali, C.~Gomez and S.~Zell,
``Quantum break-time of de Sitter,''
JCAP \textbf{06}, 028 (2017).

\bibitem{Mottola1985}
E.~Mottola,
``Particle creation in de Sitter space,''
Phys.\ Rev.\ D \textbf{31}, 754 (1985).

\bibitem{Mottola1986}
E.~Mottola,
``Thermodynamic instability of de Sitter space,''
Phys.\ Rev.\ D \textbf{33}, 1616 (1986).

\bibitem{Matarrese1998}
S.~Matarrese, S.~Mollerach and M.~Bruni,
``Second order perturbations of the Einstein-de Sitter universe,''
Phys.\ Rev.\ D \textbf{58}, 043504 (1998).

\bibitem{ACT2025}
T.~Louis et~al.\ (Atacama Cosmology Telescope Collaboration),
JCAP \textbf{11}, 062 (2025).

\bibitem{SPT3G2026}
E.~Camphuis et~al.\ (SPT-3G Collaboration),
arXiv:2506.20707.

\bibitem{Labbe2023}
I.~Labb\'e et~al.,
Nature \textbf{616}, 266 (2023).

\bibitem{GW170817}
B.~P.~Abbott et~al.\ (LIGO/Virgo),
Phys.\ Rev.\ Lett.\ \textbf{119}, 161101 (2017).

\bibitem{Alicki2023}
R.~Alicki, G.~Barenboim and A.~Jenkins,
Phys.\ Rev.\ D \textbf{108}, 123530 (2023).

\bibitem{Alicki2025}
R.~Alicki, G.~Barenboim and A.~Jenkins,
Phys.\ Lett.\ B \textbf{866}, 139519 (2025).

\bibitem{Jacobson1995}
T.~Jacobson,
``Thermodynamics of spacetime: The Einstein equation of state,''
Phys.\ Rev.\ Lett.\ \textbf{75}, 1260 (1995).

\bibitem{Scolnic2022}
D.~M.~Scolnic et~al.,
``The Pantheon+ analysis: The full data set and light-curve release,''
Astrophys.\ J.\ \textbf{938}, 113 (2022).

\bibitem{Milgrom1983}
M.~Milgrom,
``A modification of the Newtonian dynamics as a possible alternative to the hidden mass hypothesis,''
Astrophys.\ J.\ \textbf{270}, 365 (1983).


\bibitem{Planck2020}
Y.~Akrami et~al.\ (Planck Collaboration),
Astron.\ Astrophys.\ \textbf{641}, A10 (2020).

\bibitem{DESI2025}
A.~G.~Adame et~al.\ (DESI Collaboration),
JCAP \textbf{02}, 021 (2025).

\bibitem{Afshordi2017}
N.~Afshordi, C.~Coriano, L.~Delle~Rose, E.~Gould and K.~Skenderis,
Phys.\ Rev.\ Lett.\ \textbf{118}, 041301 (2017).

\end{thebibliography}

\end{document}
